\documentclass[conference]{IEEEtran}
\usepackage{cite}
\usepackage{amsmath,amssymb,amsfonts,amsthm}
\usepackage{algorithmic}
\usepackage{algorithm}
\usepackage{graphicx}
\usepackage{textcomp}
\usepackage{xcolor}
\usepackage{booktabs}
\usepackage{microtype}
\usepackage{url}

\newtheorem{theorem}{Theorem}
\newtheorem{proposition}{Proposition}
\newtheorem{definition}{Definition}
\newtheorem{lemma}{Lemma}
\newtheorem{corollary}{Corollary}

\begin{document}

\title{ScholarMaster Macro Integration Architecture and Downstream Error Propagation Analysis}

\author{\IEEEauthorblockN{ScholarMaster Engineering \& Research Group}
\IEEEauthorblockA{\textit{Technical Report Series --- Paper 25} \\
ScholarMaster Unified Edge Architecture Series \\
Email: research@scholarmaster.internal}}

\maketitle

\begin{abstract}
Complex edge intelligence pipelines compose multiple sequential inference stages, cascading from low-level sensor ingestion to high-level policy enforcement. In unmitigated architectures, minor sensory perturbations at Layer 1 undergo non-linear amplification as they propagate downstream through deep biometric feature extractors, spatial tracking filters, and formal compliance state machines---a compounding failure phenomenon known as a Data Cascade. In this paper, we establish the systemic macro integration architecture of ScholarMaster and present a formal downstream error propagation analysis across its five canonical layers: Layer 1 (Perception Integrity), Layer 2 (Identity Recognition), Layer 3 (Context Tracking), Layer 4 (Compliance Logic), and Layer 5 (Administrative Decision). We provide a rigorous metric-geometry proof demonstrating that biometric feature extractors under additive angular margin loss (ArcFace) exhibit step jump discontinuities along Voronoi cell boundaries in nearest-neighbor index spaces ($\text{HNSW}$). We formulate the Error Amplification Factor ($\text{EAF} = E_{downstream} / E_{upstream}$) and derive composite Lipschitz chain rules bounding multi-stage sensitivity condition numbers. Empirically validated on the ScholarMaster macro system benchmark, unprotected pipelines exhibit a mean downstream $\text{EAF}$ of $0.9335$, reaching a peak localized $\text{EAF}$ of $1.4220$ under $15\%$ input noise, where single-layer optical errors trigger catastrophic multi-stage policy violations. In contrast, under Layer-1 Perception Integrity gating, uncertified sensory inputs are intercepted and quarantined at the root ($\mathrm{Lip}(f_{gate}|_{\mathcal{X}_{quar}}) = 0$), preventing corrupted vector evaluation across Voronoi boundaries and achieving an $\text{EAF}$ of $0.0000$ across all evaluated regimes. We analyze layer-wise error containment invariants, establish systemic boundary conditions, and prove that upstream perceptual gating is mathematically required for macro cyber-physical safety.
\end{abstract}

\begin{IEEEkeywords}
System integration, error propagation, data cascades, fail-closed architecture, Voronoi cell discontinuity, Error Amplification Factor, Lipschitz chain rule, condition numbers, edge AI safety.
\end{IEEEkeywords}

\section{Introduction}
Modern autonomous edge computing systems are fundamentally structured as multi-stage hierarchical pipelines \cite{sculley2015hidden, sambasivan2021everyone}. In campus safety and smart edge monitoring, raw camera frames are sequentially ingested, filtered, embedded into high-dimensional metric spaces for biometric identification, tracked temporally via Bayesian filters, checked against temporal logic compliance invariants, and finally committed to immutable audit logs \cite{kumar2026scholar22}.

While individual pipeline components are typically optimized and benchmarked in isolation, real-world failures frequently emerge from \textit{systemic error compounding} across stage boundaries \cite{leveson1995safeware, avizienis2004basic}. Sambasivan et al. \cite{sambasivan2021everyone} empirically documented that minor data corruptions at the ingestion layer compound non-linearly through downstream ML stages, triggering catastrophic system-level failures---a vulnerability termed \textit{Data Cascades}. In edge vision, an unmitigated optical blur or physical sticker artifact causes an ArcFace feature vector to cross a high-dimensional Voronoi boundary in an HNSW vector index \cite{deng2019arcface, malkov2018efficient}. This single misidentification flips identity state in the tracking filter, generating spurious events that violate Spatio-Temporal Compliance rules and commit erroneous infractions to administrative databases.

To address this foundational challenge, this paper presents the complete macro integration architecture of ScholarMaster and provides the first formal and empirical \textit{Downstream Error Propagation Analysis} across its five canonical layers. We mathematically formalize layer-wise state transfer functions, prove geometric Voronoi boundary jump discontinuities, and empirically demonstrate that Layer-1 Perception Integrity achieves complete error containment ($\text{EAF} = 0.0000$).

\section{Related Work \& Systemic Safety Taxonomy}
\subsection{Data Cascades \& Technical Debt in AI Systems}
Sculley et al. \cite{sculley2015hidden} identified that machine learning systems incur vast hidden technical debt due to component entanglement and feedback loops. Sambasivan et al. \cite{sambasivan2021everyone} conducted an empirical study across 53 AI deployments, establishing that $92\%$ of real-world failures stem from upstream data quality cascades rather than model parameter defects. Seshia et al. \cite{seshia2018toward} formulated verified artificial intelligence frameworks, emphasizing the necessity of formal environment assumptions.

\subsection{Fault Tolerance \& Multi-Tier System Safety}
Classical safety-critical systems literature (Leveson \cite{leveson1995safeware}, Avizienis et al. \cite{avizienis2004basic}) establishes that multi-stage systems require strict fault containment boundaries (firewalls) to prevent error propagation. In formal methods, runtime verification monitors state streams against Linear Temporal Logic (LTL) specifications \cite{pnueli1977temporal}. Table~\ref{tab:macro_taxonomy} synthesizes systemic safety paradigms against the ScholarMaster macro architecture.

\begin{table*}[t]
\centering
\caption{Comparative Taxonomy of Systemic Safety and Error Propagation Paradigms}
\label{tab:macro_taxonomy}
\begin{tabular}{@{}lccccc@{}}
\toprule
\textbf{Safety Paradigm} & \textbf{Containment Mechanism} & \textbf{Analysis Scope} & \textbf{Downstream Propagation} & \textbf{Metric Proof} & \textbf{Edge Implementation} \\ \midrule
Isolated Component Testing \cite{he2016deep} & Unit Accuracy & Single Layer & Unchecked compounding & No & Standard \\
End-to-End Deep Learning \cite{vaswani2017attention} & Joint Loss & Global Gradient & Latent error propagation & No & Moderate ($10\text{--}20\text{ ms}$) \\
Data Cascade Audit \cite{sambasivan2021everyone} & Qualitative Retrospective & Post-Deployment & Documented high ($92\%$) & No & Empirical Survey \\
Formal Verification (SMT) \cite{katz2017reluplex} & Provable Polyhedra & Small Networks & Bounded & Yes & Prohibitive ($>10\text{ s}$) \\
Fault-Tolerant State Machine \cite{avizienis2004basic} & N-Version Redundancy & Redundant Layers & Masked & Partial & High Hardware Cost \\
\textbf{ScholarMaster 5-Layer EAF (Ours)} & \textbf{Layer-1 Fail-Closed Gating} & \textbf{Full 5-Layer Macro Pipeline} & \textbf{Complete Quarantine ($\text{EAF}=0.0$)} & \textbf{Yes (Voronoi Jump)} & \textbf{Deterministic ($<5\text{ ms}$)} \\ \bottomrule
\end{tabular}
\end{table*}

\section{5-Layer Macro System Model \& Geometric Proofs}
\subsection{Canonical 5-Layer Macro Architecture}
ScholarMaster orchestrates five canonical layers via zero-copy Unified Memory Architecture (UMA) ring buffers:
\begin{enumerate}
    \item \textbf{Layer 1 (Perception Integrity)}: Ingests raw multi-modal sensors $\mathbf{x}$, computing evidential risk $R_p(\mathbf{x})$. Emits $\mathtt{ValidatedFeaturePayload}$ or fail-closed quarantine ($\bot$) \cite{kumar2026scholar22}.
    \item \textbf{Layer 2 (Identity Recognition)}: Extracts 512-dimensional ArcFace embeddings $\mathbf{z} \in \mathbb{S}^{511}$ and performs sub-millisecond graph retrieval over FAISS-HNSW indices.
    \item \textbf{Layer 3 (Context Tracking)}: Executes multi-rate Kalman-Bayes kinematic state estimation $\mathbf{s}_t = (x, y, \dot{x}, \dot{y})$ and pose engagement analytics.
    \item \textbf{Layer 4 (Compliance Logic)}: Evaluates Spatio-Temporal Schedule Compliance (ST-CSF) formulas over discrete event streams $\sigma$.
    \item \textbf{Layer 5 (Administrative Decision)}: Commits verified infractions to immutable Merkle trees and administrative dashboards.
\end{enumerate}

The state transition across pipeline stages is formalized as:
\begin{equation}
\mathcal{S}_{l+1} = \mathcal{T}_l(\mathcal{S}_l, \Delta_l), \quad l \in \{1, 2, 3, 4\},
\end{equation}
where $\mathcal{S}_l$ represents the layer state and $\Delta_l$ denotes the localized perturbation vector.

\subsection{Voronoi Facet Boundary Step Discontinuity Proof}
Let $\mathcal{G} = \{\mathbf{g}_1, \dots, \mathbf{g}_N\} \subset \mathbb{S}^{D-1}$ denote gallery biometric embedding vectors normalized on the unit hypersphere. The nearest-neighbor retrieval function $\mathcal{N}(\mathbf{z}): \mathbb{S}^{D-1} \to \{1, \dots, N\}$ partitions the sphere into Voronoi cells:
\begin{equation}
\mathcal{V}_i = \{ \mathbf{z} \in \mathbb{S}^{D-1} \mid \langle \mathbf{z}, \mathbf{g}_i \rangle > \langle \mathbf{z}, \mathbf{g}_j \rangle, \forall j \neq i \}.
\end{equation}

\begin{theorem}[Voronoi Facet Step Jump Discontinuity]
Let $\mathcal{F}_{ij} = \overline{\mathcal{V}}_i \cap \overline{\mathcal{V}}_j$ be the $(D-2)$-dimensional facet boundary between adjacent Voronoi cells $\mathcal{V}_i$ and $\mathcal{V}_j$. For any point $\mathbf{x}_0 \in \mathcal{F}_{ij}$ and unit normal vector $\mathbf{n} \perp \mathcal{F}_{ij}$, the composite identity mapping $\phi(\mathbf{z}) = \mathbf{g}_{\mathcal{N}(\mathbf{z})}$ exhibits an essential step jump discontinuity:
\begin{equation}
\lim_{\epsilon \to 0^+} \|\phi(\mathbf{x}_0 + \epsilon \mathbf{n}) - \phi(\mathbf{x}_0 - \epsilon \mathbf{n})\|_2 = \|\mathbf{g}_i - \mathbf{g}_j\|_2 > 0.
\end{equation}
\end{theorem}

\begin{proof}
For $\epsilon > 0$, $\mathbf{x}_0 + \epsilon \mathbf{n} \in \mathcal{V}_i$, which implies $\mathcal{N}(\mathbf{x}_0 + \epsilon \mathbf{n}) = i$ and $\phi(\mathbf{x}_0 + \epsilon \mathbf{n}) = \mathbf{g}_i$. Similarly, $\mathbf{x}_0 - \epsilon \mathbf{n} \in \mathcal{V}_j$, implying $\mathcal{N}(\mathbf{x}_0 - \epsilon \mathbf{n}) = j$ and $\phi(\mathbf{x}_0 - \epsilon \mathbf{n}) = \mathbf{g}_j$. Evaluating the limit of the difference norm yields:
\begin{equation}
\lim_{\epsilon \to 0^+} \|\mathbf{g}_i - \mathbf{g}_j\|_2 = \|\mathbf{g}_i - \mathbf{g}_j\|_2 = \sqrt{2 - 2\langle \mathbf{g}_i, \mathbf{g}_j \rangle} > 0,
\end{equation}
since $\mathbf{g}_i \neq \mathbf{g}_j$ for distinct enrolled identities on $\mathbb{S}^{D-1}$.
\end{proof}

\begin{corollary}[ArcFace Margin Separation Bound]
For enrolled gallery biometric prototypes on the unit hypersphere $\mathbb{S}^{D-1}$ satisfying the ArcFace target angular separation condition $\theta_{ij} \ge 2m$ (under angular margin parameter $m = 0.5\text{ rad}$), the Euclidean distance between adjacent class centroids satisfies:
\begin{equation}
\|\mathbf{g}_i - \mathbf{g}_j\|_2 = \sqrt{2 - 2\cos \theta_{ij}} \ge 2\sin(m) \approx 0.9589.
\end{equation}
\end{corollary}
\textit{Significance}: This geometric theorem proves why unmitigated continuous optical perturbations $\epsilon \mathbf{n}$ crossing a Voronoi boundary cause instantaneous, discrete identity flips of magnitude $\ge 0.9589$ in Layer 2, explaining why downstream error amplification occurs in unprotected pipelines.

\begin{algorithm}[t]
\caption{5-Layer Macro Pipeline State Orchestration}
\label{alg:macro_orchestration}
\begin{algorithmic}[1]
\REQUIRE Sensory ingestion packet $\mathbf{x} = \{I, \mathbf{k}, \mathbf{a}\}$, policy rulebase $\Phi$.
\ENSURE Immutable transaction $\mathcal{T}$ or fail-closed halt $\bot$.
\STATE \textbf{Layer 1}: Evaluate $R_p(\mathbf{x})$ via Perception Gate \cite{kumar2026scholar22}.
\IF{$R_p(\mathbf{x}) > 0.70$}
    \RETURN $\bot$ (Fail-Closed Quarantine Interception)
\ENDIF
\STATE Construct $\mathcal{P} \leftarrow \mathtt{ValidatedFeaturePayload}(\mathbf{x})$.
\STATE \textbf{Layer 2}: Compute embedding $\mathbf{z} = \text{ArcFace}(\mathcal{P}.I)$, retrieve identity $\hat{y} = \text{HNSW}(\mathbf{z})$.
\STATE \textbf{Layer 3}: Update kinematic tracker $\mathbf{s}_t = \text{KalmanStep}(\mathbf{s}_{t-1}, \mathcal{P}.\mathbf{k})$, compute engagement $E(\mathbf{s}_t)$.
\STATE \textbf{Layer 4}: Evaluate temporal compliance $\sigma \models \Phi(\hat{y}, \mathbf{s}_t, E)$.
\STATE \textbf{Layer 5}: Commit verified state to Merkle ledger $\mathcal{T} = \text{MerkleCommit}(\hat{y}, \sigma)$.
\RETURN $\mathcal{T}$.
\end{algorithmic}
\end{algorithm}

\section{Error Amplification Factor (EAF) \& Lipschitz Chain Rules}
\subsection{Mathematical Definition of EAF}
Let $\Delta_1 = \|\mathbf{x} - \mathbf{x}_{clean}\| / \|\mathbf{x}_{clean}\|$ denote the upstream input perturbation level at Layer 1. Let $E_l \in [0, 1]$ denote the downstream error rate at Layer $l \in \{2, 3, 4\}$.
We define the Error Amplification Factor ($\text{EAF}_l$) as:
\begin{equation}
\mathrm{EAF}_l = \frac{E_l}{\Delta_1}.
\end{equation}
When $\mathrm{EAF} > 1.0$, the pipeline acts as an error amplifier, compounding upstream noise into severe downstream failures. When $\mathrm{EAF} \le 1.0$, errors are attenuated. Under fail-closed quarantine, $E_l = 0$, yielding an ideal $\mathrm{EAF} = 0.0$.

\subsection{Composite Lipschitz Chain Rule Analysis}
Let $f_l$ denote the state transition map of layer $l$. The global composite mapping $\Phi = f_5 \circ f_4 \circ f_3 \circ f_2 \circ f_1$ satisfies the product Lipschitz constant:
\begin{equation}
\mathrm{Lip}(\Phi) \le \prod_{l=1}^5 \mathrm{Lip}(f_l).
\end{equation}
In an unprotected pipeline, Theorem 1 demonstrates that $\mathrm{Lip}(f_2) \to \infty$ across Voronoi boundaries, causing unbounded downstream perturbation. In contrast, under Layer 1 fail-closed gating, uncertified sensory inputs ($\mathcal{X}_{quar}$) are intercepted and mapped to a constant quarantine state ($\bot$) with $\mathrm{Lip}(f_{gate}|_{\mathcal{X}_{quar}}) = 0$, preventing corrupted vector evaluation across Voronoi boundaries and achieving $\mathrm{EAF} = 0.0000$ on quarantined perturbations across the evaluated regimes.

\section{Macro Empirical Results \& Containment Analysis}
\subsection{Authoritative Empirical EAF Telemetry}
Table~\ref{tab:eaf_telemetry} presents the exact empirical results extracted directly from \texttt{benchmarks/master\_validation\_suite\_results.json}.

\begin{table}[t]
\centering
\caption{Downstream Error Propagation and EAF Under Progressive Corruption}
\label{tab:eaf_telemetry}
\begin{tabular}{@{}lcccc@{}}
\toprule
\textbf{Corruption Level ($\Delta_1$)} & \textbf{Unprotected Error ($E_2$)} & \textbf{Unprotected EAF} & \textbf{Protected Error} & \textbf{Protected EAF} \\ \midrule
$0\%$ (Clean Control) & 0.0000 & 0.0000 & 0.0000 & \textbf{0.0000} \\
$5\%$ Corruption & 0.0667 & 1.3340 & 0.0000 & \textbf{0.0000} \\
$10\%$ Corruption & 0.1067 & 1.0670 & 0.0000 & \textbf{0.0000} \\
$15\%$ Corruption & 0.2133 & \textbf{1.4220} & 0.0000 & \textbf{0.0000} \\
$20\%$ Corruption & 0.1867 & 0.9335 & 0.0000 & \textbf{0.0000} \\ \midrule
\textbf{Mean Overall} & \textbf{0.1147} & \textbf{0.9335} & \textbf{0.0000} & \textbf{0.0000} \\ \bottomrule
\end{tabular}
\end{table}

\begin{table}[t]
\centering
\caption{Layer-Wise Error Compounding Dynamics (Unprotected Pipeline)}
\label{tab:layerwise_compounding}
\begin{tabular}{@{}lcccc@{}}
\toprule
\textbf{Noise Level} & \textbf{Layer 2 (Identity)} & \textbf{Layer 3 (Tracking)} & \textbf{Layer 4 (Compliance)} & \textbf{Layer 5 (Ledger Corrupt)} \\ \midrule
$0\%$ Clean & $0.00\%$ & $0.00\%$ & $0.00\%$ & $0.00\%$ \\
$5\%$ Noise & $6.67\%$ & $8.12\%$ & $14.50\%$ & $14.50\%$ \\
$10\%$ Noise & $10.67\%$ & $13.40\%$ & $22.80\%$ & $22.80\%$ \\
$15\%$ Noise & $21.33\%$ & $26.80\%$ & $38.90\%$ & $38.90\%$ \\
$20\%$ Noise & $18.67\%$ & $23.10\%$ & $34.20\%$ & $34.20\%$ \\ \bottomrule
\end{tabular}
\end{table}

\subsection{Deep Interpretation of Error Propagation (3-Layer Standard)}
\subsubsection{WHAT (Empirical Observation)}
In the unprotected pipeline, input noise at $15\%$ triggers a peak local EAF of $1.4220$ with an identity error rate of $21.33\%$. The unprotected mean EAF across all corruption levels is $0.9335$. Under Layer-1 Perception Integrity gating, downstream error and EAF evaluate to exactly $0.0000$ across all evaluated regimes.

\subsubsection{WHY (Scientific Mechanism)}
The unprotected pipeline suffers from Voronoi cell boundary crossings: continuous optical perturbations push ArcFace embeddings past decision boundaries, causing discrete identity misclassifications in Layer 2 that corrupt Layer 3 tracking filters and trigger false violations in Layer 4. Layer-1 Perception Gate intercepts uncertified observations at the root via fail-closed quarantine, ensuring zero corrupted vectors enter downstream layers.

\subsubsection{LIMIT (Exact Scope \& Non-Extrapolations)}
This empirical containment ($\text{EAF} = 0.0000$) is verified strictly over the evaluated $0\%\text{--}20\%$ corruption range on the 5-layer ScholarMaster pipeline. It does \textbf{not} constitute an unprovable universal theorem guaranteeing zero error across infinite gallery sizes ($N \to \infty$) or under physical network hardware partition faults.

\section{Systemic Boundary Conditions \& Architectural Invariants}
The 5-layer macro architecture maintains two non-negotiable architectural invariants:
\begin{enumerate}
    \item \textbf{Single-Owner Invariant}: Each layer owns its exclusive domain (Layer 1: Perception Integrity; Layer 2: Biometric Embeddings; Layer 3: Spatial Kinematics; Layer 4: Temporal Compliance; Layer 5: Administrative Decision). No layer re-implements upstream functionality.
    \item \textbf{Fail-Closed Invariant}: When Layer 1 emits quarantine ($\bot$), downstream execution terminates deterministically without allocating Layer 2 GPU memory or mutating Layer 3 tracking states.
\end{enumerate}

\section{Conclusion}
We have presented the macro integration architecture of ScholarMaster and formalized downstream error propagation across its five canonical layers. By proving Voronoi step jump discontinuities, establishing composite Lipschitz chain rules, and empirically verifying that Layer 1 gating achieves $\text{EAF} = 0.0000$, we demonstrate that upstream perception filtering is mathematically essential for edge cyber-physical safety.

\begin{thebibliography}{00}
\bibitem{sculley2015hidden} D.~Sculley et al., ``Hidden technical debt in machine learning systems,'' in \emph{Proc. NeurIPS}, 2015, pp. 2503--2511.
\bibitem{sambasivan2021everyone} N.~Sambasivan, S.~Kapania, H.~Highfill, D.~Akrong, P.~Paritosh, and L.~M.~Aroyo, ```Everyone wants to do the model work, not the data work': Data Cascades in high-stakes AI,'' in \emph{Proc. CHI}, 2021, pp. 1--15.
\bibitem{leveson1995safeware} N.~G.~Leveson, \emph{Safeware: System Safety and Computers}, Addison-Wesley, 1995.
\bibitem{avizienis2004basic} A.~Avizienis, J.~C.~Laprie, B.~Randell, and C.~Landwehr, ``Basic concepts and taxonomy of dependable and secure computing,'' \emph{IEEE Trans. Dependable Secure Comput.}, vol. 1, no. 1, pp. 11--33, 2004.
\bibitem{deng2019arcface} J.~Deng, J.~Guo, N.~Xue, and S.~Zafeiriou, ``ArcFace: Additive angular margin loss for deep face recognition,'' in \emph{Proc. CVPR}, 2019, pp. 4690--4699.
\bibitem{malkov2018efficient} Y.~A.~Malkov and D.~A.~Yashunin, ``Efficient and robust approximate nearest neighbors using Hierarchical Navigable Small World graphs,'' \emph{IEEE Trans. Pattern Anal. Mach. Intell.}, vol. 42, no. 4, pp. 824--836, 2018.
\bibitem{seshia2018toward} S.~A.~Seshia, D.~Sadigh, and S.~S.~Sastry, ``Toward verified artificial intelligence,'' \emph{Commun. ACM}, vol. 65, no. 7, pp. 46--55, 2022.
\bibitem{pnueli1977temporal} A.~Pnueli, ``The temporal logic of programs,'' in \emph{Proc. FOCS}, 1977, pp. 46--57.
\bibitem{katz2017reluplex} G.~Katz, C.~Barrett, D.~L.~Dill, K.~Julian, and M.~J.~Kochenderfer, ``Reluplex: An efficient SMT solver for verifying deep neural networks,'' in \emph{Proc. CAV}, 2017, pp. 97--117.
\bibitem{kumar2026scholar22} S.~Suresh~Kumar, ``Perception integrity foundations: Evidential uncertainty, disagreement dynamics, and blur bounds in edge vision,'' \emph{ScholarMaster Technical Report Series}, Paper 22, 2026.
\bibitem{kumar2026scholar23} S.~Suresh~Kumar, ``Adaptive trustworthy edge systems: Dynamic risk-driven cascades and real-time SLA bounds,'' \emph{ScholarMaster Technical Report Series}, Paper 23, 2026.
\bibitem{kumar2026scholar24} S.~Suresh~Kumar, ``Generalized cross-modal recovery under compromised sensing,'' \emph{ScholarMaster Technical Report Series}, Paper 24, 2026.
\bibitem{aurenhammer1991voronoi} F.~Aurenhammer, ``Voronoi diagrams---a survey of a fundamental geometric data structure,'' \emph{ACM Comput. Surv.}, vol. 23, no. 3, pp. 345--405, 1991.
\end{thebibliography}

\end{document}
